% AIST / AISTATS submission draft
% Compile: pdflatex aist_submission.tex  (or xelatex for best Cyrillic in Appendix)
% Switch \documentclass to the official AISTATS style file when available.

\documentclass[11pt]{article}

\usepackage[margin=1in]{geometry}
\usepackage{amsmath,amssymb}
\usepackage{booktabs}
\usepackage{graphicx}
\usepackage{hyperref}
\usepackage{enumitem}
\usepackage{caption}
\usepackage{subcaption}
\usepackage{listings}
\usepackage{xcolor}
\usepackage[utf8]{inputenc}
\usepackage[T2A]{fontenc}
\usepackage[russian,english]{babel}

\hypersetup{colorlinks=true, linkcolor=blue, citecolor=blue, urlcolor=blue}

\lstdefinestyle{prompt}{
  basicstyle=\small\ttfamily,
  breaklines=true,
  breakatwhitespace=false,
  frame=single,
  backgroundcolor=\color{gray!5},
  columns=fullflexible,
  keepspaces=true,
}

\title{%
  Agent Cascade for Better Binary Classification on Small Datasets:\\
  Boosting Conventional ML with Vision-Language Models
}

\author{%
  Anonymous Authors\\
  \textit{(AIST submission --- anonymized draft)}
}

\date{\today}

\begin{document}
\maketitle

\begin{abstract}
Simple models such as support vector machines (SVMs) train reliably on small datasets and provide fast, interpretable segmentations, but they \emph{cannot perform perceptual analysis}: morphological cues that separate true signal from noise remain outside their representational capacity.
Conversely, vision-language models (VLMs) excel at visual reasoning yet \emph{cannot execute an SVM within their native architecture}---not without generating and running external code---and are costly as standalone classifiers on niche medical imagery.
We propose \emph{agent boosting}: a linear two-agent cascade that combines conventional ML with a general-purpose VLM, delegating only ambiguous cases to the stronger model.
A rule-based \textbf{Math Agent} screens SVM color maps using red-area coverage and spatial clustering; a \textbf{VLM Agent} (Gemini~3.1 Flash-Lite), fed few-shot exemplars and Math statistics, resolves uncertainty.
We demonstrate the approach on intraoperative optical coherence tomography (OCT): binary tumor vs.\ segmentation-noise classification on 86 scans (44 Normal, 42 Tumor).
The cascade reaches 90.5\% sensitivity, 93.2\% specificity, and 91.9\% accuracy---outperforming published SVM (83/85/84) and Gemini~3.1 Pro (86/93/90) baselines---while invoking the VLM on only 55\% of cases.
This work is an early step toward richer multi-agent graphs; code is available, the clinical dataset is withheld pending approval.
\end{abstract}

%=============================================================================
\section{Introduction}
%=============================================================================

When labeled data are scarce, practitioners still face a familiar trade-off.
\emph{Simple models}---SVMs, shallow ensembles, hand-crafted features---are fast, interpretable, and sample-efficient, yet they lack perceptual depth: they cannot distinguish a thin vertical artifact from a dense infiltrating tumor cluster without features a human would call ``visual.''
\emph{Complex models}---large VLMs---can reason about morphology, yet they cannot, within their own forward pass, train or run an SVM; doing so requires tool use, code generation, and external execution.
Replacing simple models outright wastes their strengths; using VLMs alone wastes compute and stability on easy cases.

We argue for \textbf{agent boosting}: combine conventional ML methods with VLM models so each component does what the other cannot.
A cheap deterministic agent handles segmentation-derived numerics; a general VLM is invoked only where perceptual judgment is required.
Better binary classification on small datasets emerges from \emph{composition}, not from choosing one paradigm over the other.

\paragraph{Scope.}
The system presented here is a \textbf{linear two-agent cascade}, not a full \textbf{agent graph}.
A graph formulation would require more specialized agents (e.g., morphology, calibration, cost-aware routing) and explicit edges encoding when information flows, loops, or ensembles.
We view this cascade as a deliberate first step: it already shows that boosting simple and strong models beats either alone, while remaining simple enough to analyze on $n=86$ scans.
Extending toward a multi-agent graph is ongoing work.

\paragraph{Case study.}
We instantiate agent boosting on intraoperative OCT during neurosurgery.
An SVM color-maps tissue: green (healthy white matter), yellow (damaged matter), red (tumor).
False-positive red stripes on normal tissue inflate the false-positive rate (FPR) and may mislead the surgeon.
The binary task is \textbf{TUMOR} vs.\ \textbf{NOISE} (segmentation artifact).

Our contributions:
\begin{enumerate}[nosep]
  \item Framing \emph{agent boosting} for small-data binary classification: conventional ML screening plus selective VLM escalation.
  \item A Math Agent (area, clustering, three-way verdict) and a VLM Agent (few-shot panel, structured neuro-oncology prompts).
  \item A strict cascade: VLM called only on UNCERTAIN cases (45\% of scans resolved without API).
  \item Empirical validation on 86 OCT scans vs.\ published SVM, Gemini~3.1 Pro, and HITL baselines~\cite{mdpi2024}.
\end{enumerate}

%=============================================================================
\section{Related Work}
%=============================================================================

Prior work on OCT tumor detection combines classical ML segmenters with human-in-the-loop (HITL) review by neurosurgeons~\cite{mdpi2024}.
Recent VLMs (e.g., Gemini) have been evaluated as standalone classifiers on OCT imagery.
Ensemble and stacking methods boost weak learners with stronger ones, but typically within a single model family (e.g., trees, neural nets).
Our approach differs in two ways.
First, we \emph{boost across paradigms}: geometric features from an SVM pipeline plus perceptual reasoning from a VLM.
Second, we treat routing as an \emph{agent policy}---the Math Agent decides whether the VLM is needed---anticipating richer graph-structured multi-agent systems where multiple specialists interact through typed edges.

%=============================================================================
\section{Method}
%=============================================================================

\subsection{Dataset and Task}

We evaluate agent boosting on a case study: 86 intraoperative OCT scans with SVM segmentation overlays, partitioned by histology-confirmed labels:
44 Normal (Class1, no tumor; SVM may show spurious red) and 42 Tumor (Class2, confirmed infiltration).
The binary task is: \textbf{TUMOR} vs.\ \textbf{NOISE} (artifact / false positive).
Ground truth is used only for evaluation, not for routing.

\subsection{Math Agent}

The Math Agent operates on RGB SVM maps, excluding the rightmost 15\% legend strip.
Pixels are classified by fixed color thresholds into red, yellow, green, and background (white).
Two primary features are computed over tissue pixels:

\begin{align}
  \text{red\%} &= 100 \cdot \frac{N_{\text{red}}}{N_{\text{tissue}}}, \\
  \text{clustering\%} &= 100 \cdot \frac{|C_{\max}|}{N_{\text{red}}} \quad (0 \text{ if } N_{\text{red}}=0),
\end{align}
where $|C_{\max}|$ is the largest connected component among red pixels (8-connectivity).

\textbf{Decision rules} (thresholds tuned on development scans):
\begin{itemize}[nosep]
  \item $\text{red\%} < 5\%$: NOISE, unless clustering $\geq 50\%$ and red\% $\geq 2\%$ $\Rightarrow$ UNCERTAIN.
  \item $\text{red\%} > 30\%$: TUMOR, unless clustering $< 20\%$ $\Rightarrow$ UNCERTAIN.
  \item $5\% \leq \text{red\%} \leq 30\%$: UNCERTAIN if clustering $\geq 50\%$; NOISE if clustering $< 20\%$; else UNCERTAIN.
\end{itemize}

On our dataset the Math Agent alone labels all 44 Normal scans as UNCERTAIN and 39/42 Tumor scans as TUMOR (3 Tumor scans routed to VLM).

\subsection{VLM Agent}

Model: \texttt{gemini-3.1-flash-lite}, temperature $=0$, rate-limited to 10 requests/minute.
The prompt (Appendix~\ref{app:prompts}) instructs the model as a neuro-oncology expert, provides Math Agent statistics, and embeds clinical morphology rules:
\begin{itemize}[nosep]
  \item \textbf{Noise}: thin isolated vertical red stripes surrounded by green.
  \item \textbf{Tumor}: dense red clusters infiltrating vertically toward the scan bottom, often bordering yellow damage zones.
\end{itemize}

\textbf{Few-shot panel.}
A $5\times5$ grid is rendered: cells N1--N12 (Normal exemplars, green border), T1--T12 (Tumor exemplars, red border), and TEST (center cell, thick red border).
Exemplars are stratified by red\% from the training pool (seed $=42$, manifest saved as \texttt{results/few\_shot\_exemplars.json}).
If the test scan appears among exemplars, it is replaced at panel-build time.

The model must output:
\texttt{[FINAL\_DECISION]}, \texttt{[CONFIDENCE]}, \texttt{[EXPLANATION]}.

\subsection{Cascade Policy}

\begin{center}
\begin{tabular}{lll}
\toprule
Math verdict & Action & Final label source \\
\midrule
NOISE  & Skip VLM & Math $\rightarrow$ NOISE \\
TUMOR  & Skip VLM & Math $\rightarrow$ TUMOR \\
UNCERTAIN & Call VLM & VLM $\rightarrow$ TUMOR or NOISE \\
\bottomrule
\end{tabular}
\end{center}

This \emph{strict cascade} is the simplest non-trivial agent topology: a single routing edge from Math to VLM on uncertainty.
It minimizes VLM calls while delegating perceptual disambiguation---the capability conventional ML lacks---to the stronger model.

%=============================================================================
\section{Experiments}
%=============================================================================

\subsection{Metrics}

Sensitivity (recall), specificity, and global accuracy are computed on all 86 scans with binary predictions.
Published baselines from~\cite{mdpi2024} include SVM alone, Gemini~3.1 Pro, and neurosurgeon HITL.

\subsection{Main Results}

Table~\ref{tab:metrics} summarizes performance.
The proposed cascade improves sensitivity over Gemini~3.1 Pro (+4.5~pp) and accuracy (+1.9~pp) while maintaining comparable specificity.

\begin{table}[ht]
  \centering
  \caption{Classification metrics on 86 OCT scans (\%). Baselines from~\cite{mdpi2024}.}
  \label{tab:metrics}
  \begin{tabular}{lccc}
    \toprule
    Method & Sensitivity & Specificity & Accuracy \\
    \midrule
    SVM (baseline)~\cite{mdpi2024}         & 83.0 & 85.0 & 84.0 \\
    Gemini 3.1 Pro (baseline)~\cite{mdpi2024} & 86.0 & 93.0 & 90.0 \\
    Neurosurgeons (HITL)~\cite{mdpi2024}   & 98.0 & 90.0 & 94.0 \\
    \midrule
    Math Agent only (rule-based)           & 90.5 & 88.6 & 89.5 \\
    \textbf{Math + Gemini cascade (ours)}  & \textbf{90.5} & \textbf{93.2} & \textbf{91.9} \\
    \bottomrule
  \end{tabular}
\end{table}

\textbf{Cascade routing} (86 scans): 39 decisions by Math Agent (45\%), 47 VLM calls (55\%), 0 unresolved.
Final predictions: 41 TUMOR, 45 NOISE.

\begin{table}[ht]
  \centering
  \caption{Confusion matrix for the proposed cascade.}
  \label{tab:confusion}
  \begin{tabular}{lcc}
    \toprule
    & \multicolumn{2}{c}{Predicted} \\
    \cmidrule(lr){2-3}
    Actual & NOISE & TUMOR \\
    \midrule
    Normal (44) & 41 & 3 \\
    Tumor (42)  & 4  & 38 \\
    \bottomrule
  \end{tabular}
\end{table}

The cascade recovers specificity lost by the Math Agent alone (+4.6~pp) without sacrificing sensitivity, confirming that VLM validation on UNCERTAIN Normal scans filters false positives.

\subsection{Error Analysis}

Seven misclassifications (Table~\ref{tab:errors}): 3 false positives (Normal predicted TUMOR), 4 false negatives (Tumor predicted NOISE).
All false positives occurred on VLM-routed scans (Math = UNCERTAIN); false negatives split between Math-only routing (high red\%, Math = TUMOR overridden by cascade logic---none in this run) and VLM calls on ambiguous tumor morphology.

\begin{table}[ht]
  \centering
  \caption{Misclassified scans (cascade).}
  \label{tab:errors}
  \small
  \begin{tabular}{llll}
    \toprule
    Scan ID & Final & Ground Truth & Error type \\
    \midrule
    OCT (35)Proc2 & TUMOR & Normal & FP \\
    OCT (50)Proc2 & TUMOR & Normal & FP \\
    OCT (58)Proc2 & TUMOR & Normal & FP \\
    OCT (38)Proc2 & NOISE & Tumor  & FN \\
    OCT (72)Proc2 & NOISE & Tumor  & FN \\
    OCT (76)Proc2 & NOISE & Tumor  & FN \\
    OCT (80)Proc2 & NOISE & Tumor  & FN \\
    \bottomrule
  \end{tabular}
\end{table}

\subsection{Continuous Score Analysis (LOOCV ROC)}

Treating red\% as a continuous tumor score, leave-one-out cross-validation (LOOCV) yields mean AUC $= 0.970 \pm 0.002$ with Youden-optimal threshold median $= 16.62\%$ (Table~\ref{tab:loocv}).
At this threshold, LOOCV sensitivity/specificity (90.5/88.6) match the rule-based Math Agent, indicating that the hand-crafted rules approximate an optimal univariate red\% cutoff.

\begin{table}[ht]
  \centering
  \caption{LOOCV ROC analysis using red\% as score ($n=86$ folds).}
  \label{tab:loocv}
  \begin{tabular}{lc}
    \toprule
    Quantity & Value \\
    \midrule
    AUC (mean $\pm$ std) & $0.970 \pm 0.002$ \\
    Youden $J$ & 0.791 \\
    Threshold (median across folds) & 16.62\% red \\
    Sensitivity / Specificity / Accuracy & 90.5 / 88.6 / 89.5\% \\
    TP / TN / FP / FN & 38 / 39 / 5 / 4 \\
    \bottomrule
  \end{tabular}
\end{table}

%=============================================================================
\section{Discussion}
%=============================================================================

The results support the \emph{agent boosting} hypothesis on a small medical dataset:
\begin{itemize}[nosep]
  \item \textbf{Simple alone is not enough.} The Math Agent (conventional ML features on SVM output) reaches high sensitivity but cannot auto-resolve Normal cases---all 44 are UNCERTAIN under our rules.
  \item \textbf{VLM alone is not necessary on every case.} Selective escalation restores specificity to 93.2\% with no sensitivity loss vs.\ Math-only, beating Gemini~3.1 Pro used as a standalone baseline.
  \item \textbf{Composition beats either part.} The cascade outperforms both published SVM and VLM baselines while halving API cost relative to always-on VLM inference.
\end{itemize}

Each model class contributes a capability the other lacks: SVM~+~Math for cheap numeric screening; VLM for perceptual disambiguation that SVM features cannot encode.

\paragraph{Toward agent graphs.}
The present system is a single-path cascade (Math $\rightarrow$ VLM).
A fuller agent graph would add nodes for complementary skills---morphology specialists, uncertainty calibration, ensemble voters---and edges for bidirectional feedback.
We have not yet reached that level of agent diversity or interconnection; this submission documents the boosting principle and its empirical payoff as groundwork.

\textbf{Limitations.}
Small single-center dataset ($n=86$); no confidence intervals in this draft (planned for revision).
Few-shot exemplars are drawn from the same pool as test scans (with leave-one-out replacement when a test image is selected as exemplar).
Gemini~3.1 Flash-Lite is weaker than Pro but sufficient when guided by Math features.
The linear topology may underfit settings requiring multi-hop reasoning or agent negotiation.

%=============================================================================
\section{Reproducibility}
%=============================================================================

Source code is in the project repository:
\begin{itemize}[nosep]
  \item \texttt{agents/math\_agent.py} --- Math Agent
  \item \texttt{agents/vlm\_agent.py} --- VLM Agent and prompts
  \item \texttt{agents/few\_shot.py} --- $5\times5$ panel builder
  \item \texttt{agents/pipeline.py} --- cascade orchestration
  \item \texttt{main.py} --- CLI entry point
  \item \texttt{evaluate\_roc.py} --- LOOCV ROC
\end{itemize}

\textbf{API key.}
Set \texttt{GEMINI\_API\_KEY} via environment variable or \texttt{.env} (see \texttt{.env.example}).
Do \textbf{not} commit secrets.

\textbf{Dataset.}
The \texttt{DATASETS\_MDPI/} folder is excluded from publication.
Place scans in \texttt{Class1/} (Normal) and \texttt{Class2/} (Tumor) after access approval.

\textbf{Running.}
\begin{lstlisting}[style=prompt]
pip install -r requirements.txt
cp .env.example .env   # add your API key locally
python main.py DATASETS_MDPI
python evaluate_roc.py DATASETS_MDPI
\end{lstlisting}

Full per-scan outputs: \texttt{results/scans\_summary.csv}, \texttt{results/pipeline\_results.json}.

%=============================================================================
\section{Conclusion}
%=============================================================================

We presented \emph{agent boosting} for binary classification on small datasets: conventional ML (SVM segmentation plus a Math Agent) combined with a general-purpose VLM (Gemini~3.1 Flash-Lite) in a linear two-agent cascade.
Simple models supply fast, interpretable screening; the VLM supplies perceptual analysis that SVMs cannot perform natively---and that VLMs cannot replace with internal SVM execution.
On 86 intraoperative OCT scans the cascade reaches 91.9\% accuracy, surpassing published SVM and Gemini~Pro baselines while calling the VLM on only 55\% of cases.
This is a step toward multi-agent graphs with richer specialist nodes and explicit inter-agent wiring; the boosting principle already yields measurable gains before that infrastructure is complete.

%=============================================================================
% References
%=============================================================================
\begin{thebibliography}{9}
\bibitem{mdpi2024}
Anonymous et al.
\newblock Intraoperative OCT tumor detection with SVM segmentation and LLM validation.
\newblock \emph{MDPI Journal}, 2024.
\newblock Reference [11] from the baseline study; update with full citation.
\end{thebibliography}

%=============================================================================
\appendix
%=============================================================================
\section{VLM Prompt Templates}
\label{app:prompts}

\subsection{Main prompt (Russian, as deployed)}
\begin{lstlisting}[style=prompt,basicstyle=\footnotesize\ttfamily]
У тебя роль эксперта в нейроонкологии и анализе интраоперационных
ОКТ-сканов головного мозга.
Перед тобой результаты сегментации ткани, выполненной базовым
классификатором SVM, и их математический анализ.

Цветовая кодировка на изображении:
- Зеленый = Здоровая белая ткань (Normal tissue)
- Желтый = Поврежденная ткань (Damaged matter / DM)
- Красный = Опухолевая инфильтрация (Tumor)

{few_shot_section}

ДАННЫЕ МАТЕМАТИЧЕСКОГО АГЕНТА ДЛЯ TEST-СКАНА:
- Процент площади красного цвета (опухоль): {red_percentage}%
- Процент площади желтого цвета (повреждения): {yellow_percentage}%
- Процент площади зеленого цвета (норма): {green_percentage}%
- Коэффициент кластеризации красного цвета: {clustering_ratio}%
  (высокий процент = единый очаг; низкий = рассеянный шум)
- Предварительный вердикт математического агента: {math_decision}

ТВОЯ ЗАДАЧА:
Используй эталонные примеры (N1–N12 = Normal, T1–T12 = Tumor)
как визуальную «насмотренность».
Сравни TEST-скан с эталонами по морфологии красных зон, глубине
прорастания и окружению (зелёный/жёлтый).
Прими финальное решение только для TEST-скана: TUMOR или NOISE.

Клинические правила:
1. Шум: тонкие изолированные вертикальные полосы на зелёном фоне.
2. Опухоль: плотные скопления красного, прорастающие вглубь,
   часто у жёлтой зоны.

Формат ответа:
[FINAL_DECISION]: <TUMOR или NOISE>
[CONFIDENCE]: <0%–100%>
[EXPLANATION]: <краткое обоснование>
\end{lstlisting}

\subsection{Few-shot panel section}
\begin{lstlisting}[style=prompt,basicstyle=\footnotesize\ttfamily]
FEW-SHOT ПАНЕЛЬ 5x5:
На приложенном изображении — сетка 5×5 (25 ячеек):
  • N1–N12 (зелёная рамка) — эталоны Normal
  • T1–T12 (красная рамка) — эталоны Tumor
  • TEST (толстая красная рамка, центр) — скан для классификации

Расположение ячеек (слева направо, сверху вниз):
{layout_description}
\end{lstlisting}

\subsection{English translation (for reviewers)}
\begin{lstlisting}[style=prompt,basicstyle=\footnotesize\ttfamily]
You are a neuro-oncology expert analyzing intraoperative brain OCT scans.
You receive SVM tissue segmentation (green=normal, yellow=damaged,
red=tumor) and Math Agent statistics for the TEST scan.

Use few-shot exemplars N1–N12 (Normal) and T1–T12 (Tumor) in the
attached 5x5 panel. Compare TEST morphology, infiltration depth,
and surrounding tissue. Decide TUMOR vs NOISE (segmentation artifact).

Rules: (1) Noise = thin vertical red stripes on green background;
(2) Tumor = dense red clusters infiltrating downward, often near yellow.

Output: [FINAL_DECISION], [CONFIDENCE], [EXPLANATION].
\end{lstlisting}

\section{Math Agent Thresholds}
\label{app:math}

\begin{table}[ht]
  \centering
  \caption{Math Agent constants (\texttt{agents/math\_agent.py}).}
  \begin{tabular}{ll}
    \toprule
    Parameter & Value \\
    \midrule
    Legend exclusion (right strip) & 15\% of image width \\
    Red area NOISE threshold & $< 5\%$ \\
    Red area TUMOR threshold & $> 30\%$ \\
    Clustering NOISE threshold & $< 20\%$ \\
    Clustering TUMOR threshold & $\geq 50\%$ \\
    Minimum red\% for low-area UNCERTAIN & $\geq 2\%$ \\
    \bottomrule
  \end{tabular}
\end{table}

\section{Per-Scan Summary (abbreviated)}
\label{app:scans}

Full table (86 rows) is in \texttt{results/scans\_summary.csv}.
Columns: filename, red\%, Math opinion, VLM opinion, final decision, decision source, ground truth, error flag.

\textbf{Decision source breakdown:}
Math = 39 scans (100\% of non-UNCERTAIN);
VLM = 47 scans (all Math = UNCERTAIN).

\section{Hyperparameters}
\label{app:hparams}

\begin{table}[ht]
  \centering
  \caption{Model and pipeline settings.}
  \begin{tabular}{ll}
    \toprule
    Setting & Value \\
    \midrule
    VLM model & \texttt{gemini-3.1-flash-lite} \\
    Temperature & 0.0 \\
    Max requests per minute & 10 \\
    Few-shot grid & $5 \times 5$ (12 N + 12 T + 1 TEST) \\
    Few-shot seed & 42 \\
    Exemplars per class & 12 \\
    Cell size in panel & $280 \times 210$ px \\
    Cascade & enabled (VLM only if UNCERTAIN) \\
    \bottomrule
  \end{tabular}
\end{table}

\end{document}
